\documentclass[a4paper,fleqn]{cas-dc}

\usepackage[numbers]{natbib}
\usepackage{amsmath, amsfonts, amssymb}
\usepackage{algorithm, algpseudocode}
\usepackage{graphicx}
\usepackage{adjustbox}
\usepackage{xcolor}
\usepackage{booktabs, tabularx, multirow}
\usepackage{enumitem}
\usepackage{ragged2e}
\usepackage{microtype} % Added for better typography
\setlist{nosep}

% Prevent widow and orphan lines
\widowpenalty=10000
\clubpenalty=10000

\newcommand*\D{\mathop{}\!\mathrm{d}}
\definecolor{darkblue}{RGB}{0,0,139}
\newcolumntype{B}{>{\bfseries\textcolor{darkblue}}c}
\newcolumntype{C}{>{\bfseries\textcolor{darkblue}}c}

% Stop-gradient operator for STE
\DeclareMathOperator{\sg}{sg}

\begin{document}

\shorttitle{Pricing Swing Options in Energy Markets using DRL}

\title[mode=title]{Pricing Swing Options in Energy Markets using Deep Reinforcement Learning}

\author[1]{Alexander Tsoskounoglou}[]
\ead{alexander.tsoskounoglou@gmail.com}

\author[2,3]{Sven Karbach}[]
\ead{sven@karbach.org}

\author[2,3]{Konstantinos Chatziandreou}[]
\ead{k.chatziandreou@uva.nl}

\affiliation[1]{%
  organization={Independent Researcher},
  city={Budapest},
  country={Hungary}
}

\affiliation[2]{%
  organization={Informatics Institute, University of Amsterdam},
  addressline={LAB42, Science Park 900},
  city={Amsterdam},
  postcode={1098 XH},
  country={The Netherlands}
}

\affiliation[3]{%
  organization={Korteweg-de Vries Institute for Mathematics, University of Amsterdam},
  addressline={Science Park 105–107},
  city={Amsterdam},
  postcode={1098 XG},
  country={The Netherlands}
}

% Abstract
\begin{abstract}
We present a deep reinforcement learning framework for valuing swing options in electricity markets whose spot dynamics exhibit seasonality, mean reversion, and spike risk. We formulate the swing contract as a continuous-action stochastic control problem with volume constraints, refraction periods, and convex exercise costs. The agent observes a nine-dimensional Markovian state comprising price factors, inventory status, and contract progress, and learns optimal exercise quantities using a novel D4PG algorithm enhanced with a \emph{profitability-gated actor}. This gate enforces non-negative immediate payoffs via a hard projection coupled with a straight-through estimator, enabling gradient-based learning while guaranteeing constraint satisfaction. Our approach eliminates the need for grid-based dynamic programming and density approximations, enabling scalability to high-dimensional market models and complex contractual features. Numerical results demonstrate that the learned policies significantly outperform Least-Squares Monte Carlo (LSM) benchmarks under convex exercise costs, with relative improvements exceeding 60\% in strongly convex regimes where interior optima arise.
\end{abstract}

\begin{keywords}
Deep Reinforcement Learning \sep Swing Options \sep Jump-Diffusion Models \sep Convex Transaction Costs \sep Profitability Gate \sep D4PG
\end{keywords}

\maketitle

\section{Introduction}
\label{Introduction}

Electricity markets display distinctive features such as non-storability, strong seasonality, abrupt demand-supply imbalances and short-lived price spikes. These characteristics lead to spot price dynamics that differ fundamentally from those in traditional financial markets. A large part of the literature therefore models electricity prices with mean-reverting diffusions augmented by jump components. Early two-factor Gaussian models were introduced by~\cite{lucia2002electricity}, while~\cite{cartea2005pricing} incorporated jumps to better capture extreme movements. A major refinement was provided by~\cite{Hambly2009}, who specified the log spot price as the sum of a seasonal function, an Ornstein-Uhlenbeck process and a rapidly mean-reverting jump component. This model captures both typical price movements and spike behaviour while retaining analytical tractability through transform methods.

Swing options are central to gas and power markets because they grant the holder flexible rights to purchase energy subject to daily and cumulative constraints. Their valuation is considerably more involved than that of European-style derivatives, since the holder repeatedly decides how much to exercise at each decision time. Classical valuation approaches include tree and grid-based dynamic programming~\cite{jaillet2004valuation}, Monte-Carlo regression or dual formulations~\cite{meinshausen2004monte,carmona2004swing} and analytically tractable special cases. These approaches can be accurate but face limitations when the underlying model contains multiple stochastic factors, jump components or complex contractual features. In the jump-diffusion model from~\cite{Hambly2009}, for example, swing pricing requires a two-dimensional grid over the diffusive and jump components and the transition density of the spike process must be approximated. The computational cost increases rapidly with the dimensionality of the state space, illustrating the classical curse of dimensionality.

Reinforcement learning and neuro-dynamic programming~\cite{bertsekas1996neuro,sutton1998reinforcement} provide an attractive alternative to grid-based dynamic programming. Recent advances in financial applications include model-based and model-free hedging~\cite{halperin2017qlbs,buehler2019deephedging}, optimal stopping~\cite{becker2019deepoptimalstopping} or both~\cite{tsoskounoglou2025higherorder}. These methods are well suited to swing options because the exercise decision is naturally continuous, the horizon is long and the contractual constraints are easily incorporated through the design of the state and reward.

In this paper we propose a deep reinforcement learning method for pricing swing options under jump-diffusion models with convex exercise costs. Our main contributions are:
\begin{enumerate}
    \item A \textbf{profitability-gated actor architecture} that projects proposed exercise quantities onto the feasible region ensuring non-negative immediate payoffs, while preserving gradient flow via a straight-through estimator.
    \item A comprehensive \textbf{nine-dimensional state representation} capturing spot price factors, inventory status, and contract progress under the HHK model.
    \item An \textbf{exploration strategy} combining pre-squash Gaussian noise with critic warmup, enabling effective learning in continuous action spaces with convex cost structures.
    \item Extensive \textbf{numerical experiments} demonstrating that RL significantly outperforms LSM when exercise costs are convex, with improvements exceeding 60\% under strong convexity.
\end{enumerate}

Our framework therefore connects classical stochastic modelling of electricity markets with modern reinforcement learning and provides a scalable method for valuing flexible energy contracts in environments characterised by volatility, spikes and intricate contractual design.

\clearpage

\section{Market and Contract Framework}

\subsection{Electricity Spot Price Dynamics: The HHK Model}

Electricity spot prices exhibit strong mean reversion, pronounced seasonality and sudden short-lived price spikes caused by supply shortages or unexpected outages. To capture these features we adopt the model of \cite{Hambly2009}, which specifies the logarithm of the spot price as the sum of a deterministic seasonal component, a diffusive mean-reverting factor and a rapidly mean-reverting jump component.

Let $S_t$ denote the spot price and define
\begin{align}\label{eq:model}
S_t = \exp\!\big(f(t) + X_t + Y_t\big),\end{align}
where $f(t)$ represents deterministic seasonality and the stochastic factors $X_t$ and $Y_t$ follow
\begin{align}
dX_t &= -\alpha X_t \, dt + \sigma\, dW_t, \label{eq:OU}\\
dY_t &= -\beta Y_t \, dt + J_t\, dN_t. \label{eq:jumpOU}
\end{align}
Here $W_t$ is a standard Brownian motion under the risk–neutral measure, $N_t$ is a Poisson process with constant intensity $\lambda$, and $J_t$ are iid jump sizes. Following~\cite{Hambly2009} and~\cite{cartea2005pricing}, we consider exponentially distributed jumps, $J_t \sim \mathrm{Exp}(1/\mu_J)$, which generate asymmetric and heavy-tailed spike behaviour. The parameter $\alpha>0$ governs the speed of mean reversion of the diffusive component \(X_t\), while \( \beta \gg \alpha \) ensures that spikes in \(Y_t\) decay quickly, consistent with empirical observations in power markets.

The seasonal function \(f(t)\) accounts for predictable periodic structure in electricity demand and supply and may be specified, for instance, as
\[
f(t) = \log(S_0) + A \cos(2\pi t),
\]
though more elaborate parametrisations can be used. Under this formulation, the pair $(X_t,Y_t)$ is Markovian and admits closed-form expressions for its conditional moments. In particular, the moment generating function of $ X_t + Y_t$ is available in semi-analytic form~\cite{Hambly2009}, enabling calibration to observed forward prices via the identity
\[
F_t^{[T]} = \mathbb{E}\left[S_T \mid \mathcal{F}_t \right],
\]
and the seasonal function $f(T)$ can be recovered to ensure consistency with market forwards, following the methodology of~\cite{cartea2005pricing} and~\cite{lucia2002electricity}.

The model~\eqref{eq:model} offers a balance between parsimony and realism: it captures the primary stylised features of electricity prices while remaining tractable. At the same time, its jump component makes classical grid-based dynamic programming significantly more demanding, particularly because accurate pricing requires approximations of the spike transition density. This motivates the reinforcement learning framework developed later in the paper, where the underlying dynamics can be simulated directly without requiring transition densities or discretisation of the state space.

\subsection{Swing Contract Specification}

Swing contracts are widely used in electricity and gas markets to provide flexible access to energy under uncertain demand and volatile spot prices. At each contractual decision time, the holder may choose how much energy to exercise, subject to both instantaneous bounds and cumulative volume restrictions. This flexibility makes swing contracts valuable but also challenging to price, as the optimal policy involves solving a multi-period stochastic control problem.

We consider a discrete set of exercise dates
\[
t_0 = 0 < t_1 < \dots < t_{n-1} < t_n, 
\qquad 
\Delta t = t_{i+1} - t_i,
\]
and denote by \(S_t\) the electricity spot price at time \(t\). At each decision time, the holder selects an exercise quantity
\[
q_t \in [q_{\min}, q_{\max}],
\]
where \(q_{\min}\) and \(q_{\max}\) represent contractual lower and upper bounds on the offtake rate. These bounds reflect operational and physical constraints such as minimum off-take requirements and maximum power delivery capacity.

\paragraph{Cumulative volume constraint.}
Over the full contract horizon, the total exercised quantity must not exceed a specified maximum:
\[
\sum_{t=0}^{n-1} q_t \le Q_{\max}.
\]
This cap ensures that the holder cannot exploit favourable price periods indefinitely and typically corresponds to an annual or seasonal consumption limit. The constraint introduces path dependence into the valuation problem, since remaining flexibility depends on past exercise decisions.

\paragraph{Operational constraints.}
In addition to quantity bounds and cumulative volume limitations, swing contracts may include further operational features. A commonly used restriction is a \emph{refraction period}, meaning that after exercising a strictly positive amount at time \(t\), the holder must wait a fixed number of periods before exercising again. More general constraints can also be incorporated, including ramping limits
\[
|q_{t+1} - q_t| \le R,
\]
monthly or daily minimum consumption levels, or limits on the number of exercise rights. These features expand the feasible set of admissible strategies and further complicate classical dynamic programming approaches, especially under jump–diffusion dynamics.

\paragraph{Payoff structure with convex costs.}
Let \(K\) denote the contractual strike price. Exercising quantity \(q_t\) at time \(t\) produces the payoff
\begin{equation}\label{eq:payoff}
\Pi_t(q_t)
=
q_t \, (S_t - K)^+
 - c \, q_t^{\gamma_{\mathrm{cost}}},
\end{equation}
where \(c \ge 0\) is the cost coefficient and \(\gamma_{\mathrm{cost}} \ge 1\) is the cost exponent. When \(\gamma_{\mathrm{cost}} = 1\), the cost is linear in the exercise quantity. When \(\gamma_{\mathrm{cost}} > 1\), the cost is convex, penalising large exercise quantities more heavily. This convex structure captures market frictions such as transaction fees, imbalance penalties, and operational disutility arising from large changes in consumption. If \(c = 0\), the classical no-cost swing specification is recovered.

\paragraph{Valuation objective.}
Under the risk-neutral measure, the value of the swing contract is given by
\[
V_0
=
\sup_{(q_t) \in \mathcal{Q}}
\mathbb{E}
\Bigg[
\sum_{t=0}^{n-1}
e^{-r t \Delta t}
\, \Pi_t(q_t)
\Bigg],
\]
where \(r\) denotes the risk-free interest rate, and \(\mathcal{Q}\) is the set of all exercise sequences satisfying the instantaneous bounds, cumulative volume restriction, and operational constraints. This formulation leads to a high-dimensional, continuous-control, multi-period optimisation problem that is difficult to solve using classical techniques such as tree or grid-based dynamic programming \cite{Hambly2009,jaillet2004valuation,meinshausen2004monte,carmona2004swing}. In subsequent sections we show how this structure can be naturally expressed as a Markov decision process suitable for solution by deep reinforcement learning.


\section{Benchmark: Classical Dynamic Programming}

Before introducing our reinforcement learning framework, we review the classical numerical approach for valuing swing options under the HHK model. The method, originating from \cite{Hambly2009} and closely related to the grid-based techniques of \cite{jaillet2004valuation,meinshausen2004monte,carmona2004swing}, provides an important benchmark against which we validate the performance of our deep reinforcement learning algorithm.

\subsection{State Space Discretisation and Grid Construction}

Under the HHK spot price model, the Markovian state consists of the diffusive factor \(X_t\), the spike factor \(Y_t\), the cumulative exercised volume \(Q_t\), and any auxiliary state variables modelling operational restrictions. Classical dynamic programming discretises the continuous state space by constructing structured grids:
\begin{align*}
&X_t \in \{x_1,\ldots,x_{N_X}\}, \qquad 
Y_t \in \{y_1,\ldots,y_{N_Y}\}, \qquad\\
&Q_t \in \{q_0,\ldots,q_{N_Q}\}.
\end{align*}
Even without refraction periods or ramping constraints, the two-factor HHK specification requires a two-dimensional grid for \((X,Y)\), significantly increasing complexity relative to one-factor diffusion models such as \cite{lucia2002electricity}. The grid resolution must be sufficiently fine to capture spike behaviour in the \(Y_t\) component and the mean-reverting dynamics of \(X_t\).

\subsection{Transition Probabilities Under Jump-Diffusion Dynamics}

The dynamic programming recursion requires the transition law of \((X_t,Y_t)\) over each time step \(\Delta t\). The diffusion component \(X_t\) has an exact Gaussian transition density, whereas the spike factor \(Y_t\) follows a mixture distribution arising from the Poisson jump process. As shown in \cite{Hambly2009}, the transition density of \(Y_t\) over \(\Delta t\) is approximated by truncating the Poisson distribution after \(K_{\max}\) jumps:
\[
\mathbb{P}(N_{\Delta t}=k)
= e^{-\lambda \Delta t} \frac{(\lambda \Delta t)^k}{k!}, 
\qquad k = 0,1,\ldots,K_{\max}.
\]
Conditional on the number of jumps, the distribution of \(Y_{t+\Delta t}\) admits a closed form, enabling the computation of transition probabilities to neighbouring grid points. The accuracy of this approximation depends on the grid resolution and the choice of truncation parameter \(K_{\max}\), with larger values required under high jump intensity or strong mean reversion.

\subsection{Backward Dynamic Programming Recursion}

Let \(V_t(z)\) denote the value function at time \(t\) evaluated at a grid point \(z = (x_i,y_j,q_k)\). The Bellman recursion takes the form
\begin{equation}
\begin{split}
V_t(z) = \max_{q \in \mathcal{A}(z)} \bigg[ & \Pi_t(q) + e^{-r\Delta t} \\
& \sum_{z'} \mathbb{P}(Z_{t+\Delta t}=z' \mid Z_t=z, q_t=q)\, V_{t+\Delta t}(z') \bigg],
\end{split}
\end{equation}
where the maximisation is taken over all admissible exercise quantities \(q\) that satisfy instantaneous and cumulative constraints. The recursion is initialised with the terminal condition \(V_T \equiv 0\) and evaluated backwards. This formulation requires evaluating transition expectations at each grid point and for each possible action, leading to computational cost proportional to
\[
O(n \cdot N_X \cdot N_Y \cdot N_Q \cdot |\mathcal{A}|).
\]

The original HHK implementation used grid sizes such as \(120 \times 60\) for \((X,Y)\), already requiring several minutes of computation for a single contract with daily exercise opportunities. As contractual complexity increases (e.g., through convex costs or operational constraints), the size of the required grid increases rapidly.

\subsection{Limitations of the Classical Grid Method}

Despite providing reliable benchmark valuations, the HHK grid method suffers from several structural limitations:
\begin{itemize}
    \item \textbf{Curse of dimensionality.} Each additional state variable introduces a new dimension in the grid. Extensions such as stochastic volatility, fuel-spread dynamics, forecast uncertainty, or renewable generation would render the grid-based method infeasible.
    \item \textbf{Jump dynamics.} Capturing the rapidly mean-reverting spike process requires fine grid resolution and careful handling of the truncated Poisson mixture, significantly increasing computational cost and potential numerical instability.
    \item \textbf{Continuous control.} Swing options allow continuous exercise quantities \(q_t \in [q_{\min},q_{\max}]\), which must be discretised in grid-based methods, introducing bias and limiting resolution.
    \item \textbf{Convex costs.} When costs are convex (\(\gamma_{\mathrm{cost}} > 1\)), interior optima arise, and the bang-bang structure exploited by LSM-style methods breaks down.
\end{itemize}

These limitations underline the need for alternative approaches that avoid state-space discretisation and do not rely on transition-density approximations. In Section~\ref{sec:rl_formulation}, we develop a reinforcement learning formulation that treats swing exercise as a continuous-control Markov decision process and solves it using a distributional actor--critic algorithm. This method scales naturally to high-dimensional settings and accommodates complex contract specifications, as confirmed by our numerical experiments.


\section{Reinforcement Learning Formulation}
\label{sec:rl_formulation}

We now reformulate the swing option valuation problem as a continuous-control Markov decision process (MDP) and solve it using a deep reinforcement learning algorithm. This approach avoids the state-space discretisation and transition-density approximations required in classical dynamic programming, while naturally accommodating continuous exercise quantities, convex cost structures and operational constraints. Our formulation follows the principles of neuro-dynamic programming \cite{bertsekas1996neuro,sutton1998reinforcement} and leverages recent advances in distributional actor--critic methods applied to financial decision problems \cite{halperin2017qlbs,buehler2019deephedging}.

\subsection{Markov Decision Process}
\label{sec:mdp}

The valuation problem described in Section~2.2 can be represented as a discrete-time MDP on the exercise grid $\{t_0,\dots,t_{n-1}\}$. At each time $t$, the agent observes the state
\begin{equation*}
\begin{aligned}
s_t = \Big( & \underbrace{S_t - K}_{\text{payoff}},\; \underbrace{Q^{\mathrm{used}}_t / Q_{\max}}_{\text{used frac.}},\; \underbrace{Q^{\mathrm{rem}}_t / Q_{\max}}_{\text{rem. frac.}},\; \underbrace{\tau_{\mathrm{mat}}}_{\text{time to mat.}},\; \underbrace{t/n}_{\text{progress}},\; \\
& \underbrace{S_t}_{\text{spot}},\; \underbrace{X_t}_{\text{OU}},\; \underbrace{Y_t}_{\text{jump}},\; \underbrace{\tau_{\mathrm{ref}}}_{\text{refraction}} \Big).
\end{aligned}
\end{equation*}

This nine-dimensional representation captures all information relevant to the exercise decision. Table~\ref{tab:state_components} details each component and its transformation for network input.

\begin{table}[htbp]
\centering
\caption{State representation for the swing option MDP. Components are derived from HHK dynamics and contract bookkeeping.}
\label{tab:state_components}
\begin{adjustbox}{max width=\linewidth}
\begin{tabular}{clll}
\toprule
Index & Symbol & Description & Transformation \\
\midrule
0 & $S_t - K$ & Spot minus strike & Raw (for profitability gate) \\
1 & $Q^{\mathrm{used}}_t / Q_{\max}$ & Exercised volume fraction & $[0,1]$ \\
2 & $Q^{\mathrm{rem}}_t / Q_{\max}$ & Remaining capacity fraction & $[0,1]$ \\
3 & $\tau_{\mathrm{mat}}$ & Normalised time to maturity & $[0,1]$ \\
4 & $t/n$ & Contract progress & $[0,1]$ \\
5 & $S_t$ & Raw spot price & Log-moneyness: $\log(S_t/K)$ \\
6 & $X_t$ & OU diffusive factor & Robust-scaled (median/IQR) \\
7 & $Y_t$ & Jump factor & Robust-scaled (median/IQR) \\
8 & $\tau_{\mathrm{ref}}$ & Periods since last exercise & $[0,1]$ \\
\bottomrule
\end{tabular}
\end{adjustbox}
\end{table}

The first component, $S_t - K$, is included explicitly to support the profitability gate mechanism described in Section~\ref{sec:profitability_gate}. Components 1--4 and 8 are already normalised to $[0,1]$. Component 5 is transformed to log-moneyness $\log(S_t/K)$ before network input, providing better conditioning for learning. Components 6 and 7 use robust scaling (centering by median, scaling by interquartile range) to handle the heavy-tailed jump factor $Y_t$.

The agent selects a continuous, normalised exercise quantity
\[
\tilde{q}_t \in [0, 1],
\]
which is denormalized to the contract space via
\[
q_t = q_{\min} + \tilde{q}_t \cdot (q_{\max} - q_{\min}).
\]
This determines the next state through the HHK transition dynamics
\[
(X_{t+\Delta t}, Y_{t+\Delta t}) 
= 
\text{SimHHK}(X_t, Y_t; \Delta t),
\qquad
Q^{\mathrm{used}}_{t+\Delta t}
=
Q^{\mathrm{used}}_t + q_t,
\]
where $\text{SimHHK}$ denotes an Euler–Maruyama simulation of the stochastic differential equations \eqref{eq:OU}--\eqref{eq:jumpOU}.

The reward for exercising quantity $q_t$ at time step $t$ is defined as the discounted incremental payoff:
\[
r_t(q_t)
=
\delta^t
\Big(
q_t \, (S_t - K)^+
 - c \, q_t^{\gamma_{\mathrm{cost}}}
\Big),
\]
where $\delta = e^{-r\Delta t}$ is the per-step discount factor. Note that we use $\delta$ for the RL discount factor to avoid confusion with $\gamma_{\mathrm{cost}}$, the cost exponent. The objective of the agent is to maximise the expected cumulative reward
\[
\mathbb{E}\Bigg[
\sum_{t=0}^{n-1} r_t(q_t)
\Bigg],
\]
which coincides with the swing option value $V_0$ under the risk-neutral measure.

\subsection{Enforcement of Swing Constraints}

Swing contracts impose restrictions on the path of admissible actions, including instantaneous bounds, cumulative volume limits and operational rules such as refraction periods. While such constraints greatly complicate grid-based dynamic programming, they can be enforced efficiently in reinforcement learning.

\paragraph{Action clipping.}
At each decision time, the actor network outputs a normalised action $\tilde{q}_t \in [0,1]$, which is mapped to the admissible interval $[q_{\min},q_{\max}]$. This ensures instantaneous feasibility.

\paragraph{Cumulative volume bookkeeping.}
The state variable $Q^{\mathrm{used}}_t$ records exercised volume. If $Q^{\mathrm{used}}_t + q_t > Q_{\max}$, the action is clipped to the maximum feasible value, i.e.\ $q_t = Q_{\max} - Q^{\mathrm{used}}_t$.

\paragraph{Refraction period masking.}
If the contract prohibits exercise for $\tau_{\mathrm{ref}}$ periods following a non-zero action, we augment the state with a counter $\ell_t$ and mask all positive actions during the refraction window, forcing $q_t = 0$.

\paragraph{Profitability masking.}
If the proposed action would result in negative immediate payoff ($\Pi_t(q_t) < 0$), the action is masked to $q_t = 0$. This is in addition to the profitability gate described below, providing a safety net during exploration.


\subsection{D4PG Algorithm: Core Architecture}
\label{sec:d4pg}

We employ a variant of the Deterministic Policy Gradient algorithm~\cite{silver2014dpg}, enhanced with distributional value estimation and prioritized experience replay as in D4PG~\cite{barth2018d4pg}. This section describes each component and its role in swing option valuation.

\subsubsection{Deterministic Policy Gradient}
The actor network $\pi_\theta: \mathcal{S} \to [0,1]$ maps states to normalised exercise quantities. The policy is trained to maximise the critic's value estimate:
\[
\nabla_\theta J(\theta) = \mathbb{E}_{s \sim \mathcal{D}}\left[\nabla_a Q_\phi(s,a)\big|_{a=\pi_\theta(s)} \nabla_\theta \pi_\theta(s)\right].
\]
Stability is ensured through target networks $(\pi_{\theta'},Q_{\phi'})$ updated via Polyak averaging:
\[
\theta' \leftarrow (1-\tau)\theta' + \tau \theta, \quad \phi' \leftarrow (1-\tau)\phi' + \tau \phi,
\]
with $\tau = 0.0032$.

\subsubsection{Critic Architecture}
The critic $Q_\phi: \mathcal{S} \times \mathcal{A} \to \mathbb{R}$ estimates the expected cumulative discounted reward. For swing options with convex costs, the critic must capture the complex dependence of continuation value on the current exercise decision. The critic is trained to minimise the temporal difference loss:
\begin{equation*}
\begin{split}
\mathcal{L}_{\text{critic}} = \mathbb{E}_{(s,a,r,s',d) \sim \mathcal{D}}\Big[ w_i \big( & r + \delta (1-d) Q_{\phi'}(s', \pi_{\theta'}(s')) \\
& - Q_\phi(s,a) \big)^2 \Big],
\end{split}
\end{equation*}
where $w_i$ are importance sampling weights from prioritised replay and $d$ is the terminal indicator.

\subsubsection{Prioritised Experience Replay}
We employ soft prioritised experience replay with a priority schedule designed to balance exploration breadth with learning focus. Transition $i$ is sampled with probability:
\[
P(i) = \frac{p_i^\alpha}{\sum_k p_k^\alpha},
\]
where $p_i = |\delta_i| + \epsilon$ is the TD-error-based priority and $\epsilon = 5 \times 10^{-6}$ is the priority floor.

The priority exponent~$\alpha$ is scheduled: starting at $\alpha_0 = 0.1$ (near-uniform) and ramping to $\alpha_{\text{final}} = 0.2$ between episodes 5{,}000 and 25{,}000. Importance sampling weights $w_i = (N \cdot P(i))^{-\beta}$ with $\beta = 1.0$ provide full correction for the non-uniform sampling, reducing bias while maintaining sample efficiency.


\subsection{Profitability-Gated Actor}
\label{sec:profitability_gate}

A key challenge in swing option pricing with convex exercise costs is ensuring that the learned policy never proposes actions that would result in immediate losses. Under the payoff structure~\eqref{eq:payoff}
\[
\Pi(q) = q \cdot (S_t - K)^+ - c \cdot q^{\gamma_{\mathrm{cost}}},
\]
profitability depends non-trivially on the exercise quantity $q$ when $\gamma_{\mathrm{cost}} > 1$.

\subsubsection{Break-Even Quantity}
The break-even quantity $q^*$ at which revenue equals cost satisfies:
\[
q^* \cdot (S_t - K)^+ = c \cdot (q^*)^{\gamma_{\mathrm{cost}}}.
\]
For $\gamma_{\mathrm{cost}} > 1$, this yields the closed-form solution:
\[
q^* = \left(\frac{(S_t - K)^+}{c \cdot \gamma_{\mathrm{cost}}}\right)^{\frac{1}{\gamma_{\mathrm{cost}} - 1}}.
\]
Any $q > q^*$ would result in negative immediate profit.

\subsubsection{Hard Gate Mechanism}
We enforce profitability by projecting the actor's output onto the feasible region. Let $\tilde{q} \in [0,1]$ denote the raw (normalised) actor output. The gated action is:
\[
q_{\text{gated}} = \min\left(\tilde{q}, \frac{q^* - q_{\min}}{q_{\max} - q_{\min}}\right),
\]
ensuring the denormalised quantity $q_{\text{actual}} = q_{\min} + q_{\text{gated}} \cdot (q_{\max} - q_{\min})$ satisfies $\Pi(q_{\text{actual}}) \ge 0$.

\subsubsection{Straight-Through Estimator}
To enable gradient-based learning despite the non-differentiable projection, we employ a straight-through estimator (STE)~\cite{bengio2013ste}:
\[
q_{\text{out}} = \tilde{q} + \sg(q_{\text{gated}} - \tilde{q}),
\]
where $\sg(\cdot)$ denotes the stop-gradient operator. In the forward pass, $q_{\text{out}} = q_{\text{gated}}$; in the backward pass, gradients flow through $\tilde{q}$ as if no projection occurred. This approach preserves the actor's ability to learn while guaranteeing constraint satisfaction.

\subsubsection{Implementation in Actor Network}
Algorithm~\ref{alg:profitability_gate} summarises the gated actor forward pass.

\begin{algorithm}[htbp]
\caption{Profitability-Gated Actor Forward Pass}
\label{alg:profitability_gate}
\begin{algorithmic}[1]
\Require State $s_t$ with $(S_t - K)$ at index 0
\Require Cost parameters $c, \gamma_{\mathrm{cost}}, q_{\min}, q_{\max}$
\Ensure Feasible normalised action $q_{\text{out}} \in [0,1]$
\State $u \gets \text{MLP}(s_t)$ \Comment{Pre-activation output}
\State $\tilde{q} \gets \sigma(\beta \cdot u)$ \Comment{$\beta$-sigmoid with $\beta=3.0$}
\If{$c > 0$ and $\gamma_{\mathrm{cost}} > 1$}
    \State $q^* \gets \left(\frac{(S_t-K)^+}{c \cdot \gamma_{\mathrm{cost}}}\right)^{1/(\gamma_{\mathrm{cost}}-1)}$
    \State $q_{\text{limit}} \gets \max\left(0, \frac{q^* - q_{\min}}{q_{\max} - q_{\min}}\right)$
    \State $q_{\text{gated}} \gets \min(\tilde{q}, q_{\text{limit}})$
\Else
    \State $q_{\text{gated}} \gets \tilde{q}$
\EndIf
\State \Return $\tilde{q} + \sg(q_{\text{gated}} - \tilde{q})$ \Comment{STE}
\end{algorithmic}
\end{algorithm}


\subsection{Exploration Strategy}
\label{sec:exploration}

Effective exploration is critical for learning in continuous action spaces. We employ several complementary mechanisms.

\subsubsection{Pre-Squash Gaussian Noise}
Unlike standard DDPG which adds noise to the output action, we add noise to the pre-activation values before the output nonlinearity:
\[
u' = u + \sigma \cdot \epsilon, \quad \epsilon \sim \mathcal{N}(0, 1).
\]
This ``pre-squash'' noise maintains exploration even when the policy saturates near the action bounds.

The noise standard deviation follows a plateau-then-decay schedule:
\[
\sigma(e) = \begin{cases}
\sigma_0 & e < e_{\text{plateau}} \\
\sigma_{\text{floor}} + \dfrac{\sigma_0 - \sigma_{\text{floor}}}{1 + (e - e_{\text{plateau}})/e_{\text{plateau}}} & \text{otherwise}
\end{cases}
\]
with $\sigma_0 = 1.30$, $\sigma_{\text{floor}} = 0.26$, and $e_{\text{plateau}} = 3{,}200$ episodes.

Additionally, we apply adaptive scaling: $\sigma_{\text{scaled}} = \sigma \cdot (1 + 0.6 \cdot |u|)$, which increases noise when the pre-activation magnitude is large, preventing saturation lock-in.

\subsubsection{Critic Warmup Phase}
During the first 1{,}024 episodes, actor updates are frozen while only the critic learns. This allows the Q-function to stabilise before the policy begins adapting. During warmup:
\begin{itemize}
    \item Exploration noise is reduced to 40\% of its normal level
    \item Noise fraction ramps linearly from 40\% to 100\% over the warmup period
\end{itemize}

\subsubsection{Target Policy Smoothing}
To reduce overestimation in the target Q-values, we add noise to the target policy:
\[
\tilde{a}' = \pi_{\theta'}(s') + \epsilon_{\text{target}}, \quad \epsilon_{\text{target}} \sim \mathcal{N}(0, \sigma_{\text{target}}^2).
\]
The target noise starts at $\sigma_{\text{target}} = 0.15$ and decays linearly to $0.04$ after episode 20{,}000.


\subsection{Neural Network Architecture}
\label{sec:architecture}

Both actor and critic networks are lightweight multilayer perceptrons designed for sample efficiency. Table~\ref{tab:network_arch} summarises the architecture.

\begin{table}[htbp]
\centering
\caption{Neural network architecture for actor and critic.}
\label{tab:network_arch}
\begin{adjustbox}{max width=\linewidth}
\begin{tabular}{lll}
\toprule
Component & Actor & Critic \\
\midrule
Input dimension & 9 (state) & 9 (state) + 1 (action) \\
Hidden layers & 2 & 2 \\
Hidden units & 64 & 64 \\
Activation & SiLU & SiLU \\
Normalisation & LayerNorm & LayerNorm \\
Output dimension & 1 & 1 \\
Output activation & $\beta$-sigmoid ($\beta=3.0$) & None \\
\bottomrule
\end{tabular}
\end{adjustbox}
\end{table}

\paragraph{Initialisation.}
Weights are initialised using orthogonal initialisation with gain $\sqrt{2}$ for hidden layers (appropriate for ReLU-family activations). Output layer weights use uniform initialisation in $[-0.003, 0.003]$.

\paragraph{Input Preprocessing.}
The raw state undergoes domain-specific transformations before entering the network:
\begin{itemize}
    \item $S_t$ is transformed to log-moneyness: $\log(S_t / K)$
    \item $X_t$ and $Y_t$ are robust-scaled using median and interquartile range
    \item Other components pass through unchanged
\end{itemize}

\paragraph{Output Activation.}
The actor uses a $\beta$-sigmoid output: $\sigma(\beta \cdot u)$ with $\beta = 3.0$. This provides softer saturation gradients compared to standard sigmoid, improving learning in boundary regions.

\paragraph{Optimiser Configuration.}
Training uses AdamW with separate learning rates and weight decay:
\begin{itemize}
    \item Actor: $\eta_a = 1.6 \times 10^{-4}$, weight decay $5 \times 10^{-5}$
    \item Critic: $\eta_c = 9.0 \times 10^{-5}$, weight decay $1.2 \times 10^{-4}$
\end{itemize}
Learning rates follow a warmup (1{,}024 episodes) followed by cosine decay to 20\% of initial values over 40{,}000 episodes.


\subsection{Training Procedure}
\label{sec:training}

Algorithm~\ref{alg:training} presents the complete training procedure for swing option pricing.

\begin{algorithm}[htbp]
\caption{D4PG Training for Swing Option Pricing}
\label{alg:training}
\begin{algorithmic}[1]
\Require Swing contract $\mathcal{C}$, HHK parameters $\Theta_{\text{HHK}}$
\Require Training episodes $N_{\text{train}}$, evaluation paths $N_{\text{eval}}$
\Ensure Trained actor $\pi_\theta$
\State Generate paths: $(t, S, X, Y) \gets \textsc{SimHHK}(\Theta_{\text{HHK}}, N_{\text{train}} + N_{\text{eval}})$
\State Initialise actor $\pi_\theta$, critic $Q_\phi$, targets $\pi_{\theta'}, Q_{\phi'}$
\State Initialise replay buffer $\mathcal{D}$ with capacity $2 \times 10^5$
\State $\pi_\theta.\textsc{SetProfitabilityParams}(q_{\min}, q_{\max}, c, \gamma_{\text{cost}})$
\For{episode $e = 1$ to $N_{\text{train}}$}
    \State $s_0 \gets \textsc{Env.Reset}(\text{path\_idx} = e)$
    \For{step $t = 0$ to $n-1$}
        \State $\sigma \gets \textsc{NoiseSchedule}(e)$
        \If{$e \le e_{\text{warmup}}$}
            \State $\sigma \gets \sigma \cdot (0.4 + 0.6 \cdot e / e_{\text{warmup}})$ \Comment{Gradual warmup noise}
        \EndIf
        \State $a_t \gets \pi_\theta(s_t) + \sigma \cdot \epsilon$, $\epsilon \sim \mathcal{N}(0,1)$
        \State $s_{t+1}, r_t, \text{done} \gets \textsc{Env.Step}(a_t)$
        \State $\mathcal{D}.\textsc{Add}(s_t, a_t, r_t, s_{t+1}, \text{done})$
        \If{$|\mathcal{D}| \ge N_{\min}$ and $t \mod 2 = 0$}
            \State $(s, a, r, s', d, w) \gets \mathcal{D}.\textsc{Sample}(128)$
            \State Update $Q_\phi$ via TD loss with IS weights $w$
            \If{$e > e_{\text{warmup}}$}
                \State Update $\pi_\theta$ via policy gradient
                \State Soft-update $\pi_{\theta'}$
            \EndIf
            \State Soft-update $Q_{\phi'}$
        \EndIf
    \EndFor
    \If{$e \mod 1024 = 0$}
        \State Evaluate and log RL/LSM prices
    \EndIf
\EndFor
\end{algorithmic}
\end{algorithm}

\paragraph{Hyperparameter Configuration.}
Table~\ref{tab:hyperparams} lists the complete hyperparameter configuration used in experiments.

\begin{table}[htbp]
\centering
\caption{Training hyperparameters for the convex cost experiments.}
\label{tab:hyperparams}
\begin{adjustbox}{max width=\linewidth}
\begin{tabular}{llr}
\toprule
Category & Parameter & Value \\
\midrule
\multirow{4}{*}{Training} & Episodes & 32{,}768 \\
 & Batch size & 128 \\
 & Replay buffer size & 200{,}000 \\
 & Minimum replay size & 18{,}000 \\
\midrule
\multirow{3}{*}{Target Networks} & Polyak $\tau$ & 0.0032 \\
 & Target noise $\sigma$ & 0.15 $\to$ 0.04 \\
 & Target noise clip & 0.25 \\
\midrule
\multirow{4}{*}{PER} & Initial $\alpha$ & 0.1 \\
 & Final $\alpha$ & 0.2 \\
 & $\beta$ & 1.0 (constant) \\
 & Priority floor & $5 \times 10^{-6}$ \\
\midrule
Critic Warmup & Episodes & 1{,}024 \\
\bottomrule
\end{tabular}
\end{adjustbox}
\end{table}


\subsection{LSM Benchmark Methodology}
\label{sec:lsm}

To validate the RL valuations, we compare against Least-Squares Monte Carlo (LSM) following~\cite{longstaff2001valuing}. We employ an out-of-sample evaluation protocol to ensure fair comparison.

\subsubsection{Estimator Training}
LSM continuation values are estimated by regressing discounted future rewards onto polynomial basis functions. We use Chebyshev polynomials $T_k$ of degree 7, which provide numerical stability over a wide price range. The price is first normalised to $[-1, 1]$:
\[
x = \frac{2S_t - S_{\min} - S_{\max}}{S_{\max} - S_{\min}},
\]
and the regression features are:
\[
\Phi(x) = \left[T_0(x), T_1(x), \ldots, T_7(x)\right].
\]

\subsubsection{Out-of-Sample Evaluation}
We train the LSM estimators on the same 32{,}768 paths used for RL training, then evaluate both methods on a separate set of 65{,}536 held-out paths. This prevents look-ahead bias and ensures comparable evaluation conditions.

\subsubsection{Exercise Decision}
At each decision time $t$, LSM exercises quantity $q_{\max}$ if:
\[
\Pi_t(q_{\max}) + \hat{C}_t^{\text{ex}}(S_t) > \hat{C}_t^{\text{keep}}(S_t),
\]
where $\hat{C}_t^{\text{ex}}$ and $\hat{C}_t^{\text{keep}}$ are the fitted continuation values after exercising or holding, respectively.

Note that LSM is inherently designed for discrete exercise decisions (exercise $q_{\max}$ or 0), making it suboptimal for the continuous exercise problem with convex costs where interior optima exist.

\textbf{Implementation caveat.} When convex costs are present ($c > 0$, $\gamma_{\mathrm{cost}} > 1$), the terminal time step must gate exercises on \emph{net profitability} ($\Pi_T(q) > 0$) rather than merely in-the-money status ($q \cdot (S_T - K)^+ > 0$). Omitting the cost from this check leads to unprofitable exercises at the final step, biasing the LSM price downward.


\section{Numerical Experiments}
\label{sec:experiments}

\subsection{Experimental Setup}

We evaluate the proposed D4PG framework across a grid of convex cost parameters. Table~\ref{tab:contract_params} summarises the swing contract and HHK model configuration.

\begin{table}[htbp]
\centering
\caption{Contract and model parameters for numerical experiments.}
\label{tab:contract_params}
\begin{adjustbox}{max width=\linewidth}
\begin{tabular}{lllr}
\toprule
Category & Parameter & Symbol & Value \\
\midrule
\multirow{6}{*}{Contract} & Strike & $K$ & 1.0 \\
 & Maturity & $T$ & 0.0833 years \\
 & Decision dates & $n$ & 22 \\
 & Max exercise/period & $q_{\max}$ & 2.0 \\
 & Max total volume & $Q_{\max}$ & 20.0 \\
 & Risk-free rate & $r$ & 5\% \\
\midrule
\multirow{6}{*}{HHK} & Initial spot & $S_0$ & 1.0 \\
 & Mean-reversion & $\alpha$ & 12.0 \\
 & Diffusion vol. & $\sigma$ & 1.2 \\
 & Jump decay & $\beta$ & 150.0 \\
 & Jump intensity & $\lambda$ & 6.0 \\
 & Mean jump size & $\mu_J$ & 0.3 \\
\midrule
\multirow{2}{*}{Costs} & Coefficient & $c$ & \{0.01, 0.02, 0.04, 0.05, 0.08, 0.10, 0.15\} \\
 & Exponent & $\gamma_{\text{cost}}$ & \{1.0, 1.5, 2.0, 3.0\} \\
\bottomrule
\end{tabular}
\end{adjustbox}
\end{table}

Each configuration is trained with three seeds $\{11, 12, 13\}$. We report the best seed's result at the episode achieving highest RL price on evaluation paths.

\subsection{Results}

Table~\ref{tab:results} presents the comparative valuation of swing options using D4PG and LSM benchmarks. We report 95\% confidence intervals for the relative improvement ($\Delta\%$) using the delta method to approximate the variance of the price ratio. Standard errors for the individual RL and LSM prices are consistently between 0.01 and 0.02 across all configurations.

\begin{table}[htbp]
\centering
\caption{Swing option prices under convex exercise costs. $\Delta\%$ shows relative improvement of RL over LSM, with 95\% confidence intervals. Standard errors for individual prices are consistently $\approx 0.02$.}
\label{tab:results}
\begin{adjustbox}{max width=\linewidth}
\begin{tabular}{ccrrrc}
\toprule
$c$ & $\gamma$ & LSM Price & RL Price & $\Delta\%$ ($\pm 95\% \text{CI}$) \\
\midrule
0.00 & 1.0 & 2.6612 & 2.6585 & $-0.10 \pm 0.10$ \\
0.01 & 1.0 & 2.5478 & 2.5411 & $-0.26 \pm 0.09$ \\
0.01 & 1.5 & 2.5026 & 2.4992 & $-0.14 \pm 0.09$ \\
0.01 & 2.0 & 2.4402 & 2.4472 & $+0.29 \pm 0.08$ \\
0.01 & 3.0 & 2.2330 & 2.3164 & $+3.73 \pm 0.11$ \\
\midrule
0.02 & 1.0 & 2.4402 & 2.4318 & $-0.34 \pm 0.08$ \\
0.02 & 1.5 & 2.3526 & 2.3608 & $+0.35 \pm 0.07$ \\
0.02 & 2.0 & 2.2350 & 2.2692 & $+1.53 \pm 0.10$ \\
0.02 & 3.0 & 1.8593 & 2.0880 & $+12.30 \pm 0.18$ \\
\midrule
0.04 & 1.0 & 2.2302 & 2.2248 & $-0.24 \pm 0.12$ \\
0.04 & 1.5 & 2.0687 & 2.0963 & $+1.33 \pm 0.11$ \\
0.04 & 2.0 & 1.8622 & 1.9735 & $+5.98 \pm 0.14$ \\
0.04 & 3.0 & 1.2595 & 1.7957 & $+42.57 \pm 0.53$ \\
\midrule
0.05 & 1.0 & 2.1343 & 2.1302 & $-0.19 \pm 0.08$ \\
0.05 & 1.5 & 1.9327 & 1.9757 & $+2.22 \pm 0.12$ \\
0.05 & 2.0 & 1.6960 & 1.8471 & $+8.91 \pm 0.14$ \\
0.05 & 3.0 & 1.0342 & 1.6932 & $+63.72 \pm 0.84$ \\
\midrule
0.08 & 1.0 & 1.8622 & 1.8519 & $-0.55 \pm 0.08$ \\
0.08 & 1.5 & 1.5928 & 1.6568 & $+4.02 \pm 0.13$ \\
0.08 & 2.0 & 1.2277 & 1.5391 & $+25.36 \pm 0.30$ \\
\midrule
0.10 & 1.0 & 1.6960 & 1.6938 & $-0.13 \pm 0.08$ \\
0.10 & 1.5 & 1.3892 & 1.4857 & $+6.95 \pm 0.13$ \\
0.10 & 2.0 & 1.0380 & 1.3775 & $+32.71 \pm 0.45$ \\
\midrule
0.15 & 1.0 & 1.2958 & 1.3247 & $+2.23 \pm 0.10$ \\
0.15 & 1.5 & 0.9696 & 1.1009 & $+13.54 \pm 0.27$ \\
\bottomrule
\end{tabular}
\end{adjustbox}
\end{table}

\subsection{Analysis}

Several patterns emerge from Table~\ref{tab:results}:

\begin{enumerate}
    \item \textbf{Linear costs ($\gamma = 1$):} RL and LSM produce similar valuations, as expected. Both identify the bang-bang optimal policy (exercise $q_{\max}$ or 0). Minor differences arise from policy approximation.
    
    \item \textbf{Moderate convexity ($\gamma = 1.5, 2.0$):} RL increasingly outperforms LSM. The polynomial basis in LSM struggles to capture the smooth interior optima that arise with convex costs. RL improvements range from 1\% to 25\%.
    
    \item \textbf{Strong convexity ($\gamma = 3.0$):} RL achieves substantial improvements (up to 64\% for $c=0.05$, $\gamma=3.0$). At high convexity, optimal exercise quantities lie well inside $[0, q_{\max}]$, and LSM's discrete exercise assumption leads to significant undervaluation.
    
    \item \textbf{Cost coefficient dependence:} Higher cost coefficients amplify the advantage of RL, as the region of non-trivial interior optima expands.
\end{enumerate}

The profitability gate (Section~\ref{sec:profitability_gate}) is essential for learning effective policies under strong convexity, as it prevents exploration of clearly suboptimal high-exercise actions that would incur losses.

\subsection{Computational Efficiency}

Training for 32{,}768 episodes requires approximately 2--3 hours on a 4-core CPU. Once trained, the policy delivers instantaneous valuations: evaluating 65{,}536 paths takes less than 10 seconds. In contrast, LSM must be re-run for each new set of paths and does not amortise across evaluations.


\section{Conclusion}

We have developed a deep reinforcement learning framework for pricing swing options with convex exercise costs under the HHK jump-diffusion model. Our main innovation is the profitability-gated actor, which employs a hard projection to ensure non-negative immediate payoffs while preserving gradient flow through a straight-through estimator.

Numerical experiments demonstrate that the proposed method matches LSM performance under linear costs and significantly outperforms it under convex costs, with improvements exceeding 60\% in strongly convex regimes. The framework naturally accommodates high-dimensional state representations, complex operational constraints, and continuous exercise decisions without requiring state-space discretisation or transition density approximations.

Future work includes extension to multi-factor stochastic volatility models, incorporation of renewable generation uncertainty, and application to gas storage and other flexible energy contracts.


\section*{Acknowledgement}

The authors gratefully acknowledge financial support from the Amsterdam University Fund through the AUF Impact Call (Fall 2024). Additional support from the AI4FinTech initiative at the University of Amsterdam is also gratefully acknowledged.

% References
\bibliographystyle{cas-model2-names}
\bibliography{Bibliography.bib}


\end{document}

